% Capítulo textual do trabalho.

\chapter{Alíneas}
\label{cap:alineas}

Este capítulo apresenta o modo de uso de alíneas. Em documentos acadêmicos, o uso correto de alíneas é fundamental para garantir a clareza na apresentação de informações que não possuem um título próprio. A norma técnica da Associação Brasileira de Normas Técnicas (ABNT) orienta a estruturação adequada das alíneas de forma a facilitar a compreensão e a organização textual. O uso das alíneas segue diretrizes específicas que devem ser rigorosamente aplicadas, conforme descrito a seguir~\cite{abnt2012}.

As alíneas são empregadas quando é necessário enumerar vários assuntos em uma seção sem título. Um exemplo de alíneas pode ser observado na classificação dos paradigmas de programação. Os principais paradigmas de programação são:

\begin{alineas}
    \item Programação imperativa;
    \item Programação funcional;
    \item Programação orientada a objetos;
    \item Programação lógica:
        \begin{subalineas}
            \item baseada no uso de lógica formal para expressar programas;
            \item foca em resolver problemas declarando relações e permitindo que o computador derive soluções;
            \item exemplos comuns incluem linguagens como Prolog.
        \end{subalineas}
\end{alineas}

Neste exemplo, observa-se o uso correto de alíneas para organizar a informação de maneira hierárquica e clara. O uso de subalíneas foi aplicado para detalhar a ``Programação lógica''. Este tipo de estruturação facilita a compreensão e o destaque de elementos dentro de uma lista de assuntos.
